% ECONOMICS_PROBLEM: {"id": "H26-635", "title": "Third-party reporting versus privacy and compliance cost", "jel_code": "H26", "source_id": "OP-0049"}
% !TeX program = xelatex
\documentclass[11pt,a4paper]{article}
\usepackage[margin=23mm]{geometry}
\usepackage{fontspec}
\setmainfont{Latin Modern Roman}
\usepackage{amsmath,amssymb,mathtools}
\usepackage{xcolor,enumitem,booktabs,longtable,array,xurl,hyperref}
\definecolor{muted}{HTML}{53636C}
\hypersetup{colorlinks=true,urlcolor=blue,linkcolor=blue}
\setlength{\parindent}{0pt}
\setlength{\parskip}{5pt}
\newcommand{\R}{\mathbb R}
\newcommand{\E}{\mathbb E}
\newcommand{\N}{\mathbb N}
\newcommand{\ind}{\mathbf 1}
\newcommand{\Var}{\operatorname{Var}}
\newcommand{\Cov}{\operatorname{Cov}}
\newcommand{\supp}{\operatorname{supp}}
\DeclareMathOperator*{\argmin}{arg\,min}
\DeclareMathOperator*{\argmax}{arg\,max}
\newcommand{\entrymeta}[3]{{\sffamily\small\color{muted}\textbf{\texttt{#1}}\quad|\quad #2\par #3\par}}
\newcommand{\Problemref}[1]{\texttt{#1}}
\newcommand{\JELref}[2]{\texttt{#2} (JEL \texttt{#1})}
\begin{document}
\section*{Third-party reporting versus privacy and compliance cost}
\textbf{Problem ID:} \texttt{H26-635}\quad \textbf{JEL:} \texttt{H26}\par
% BEGIN ECONOMICS STATEMENT
\label{problem:OP-0049}
\label{atlas:prob:P9}
\noindent\textbf{Original atlas identifier:} P9; entry 49. \textbf{Field:} Public finance and taxation. \textbf{Original tractability label:} Medium.

\noindent\textbf{Classification:} Parameterized theoretical/empirical research agenda; not a certified unsolved theorem.

\subsection*{Definitions and assumptions}
Let reporting scope $d\in\mathcal D$ specify data fields, transaction coverage, retention, verification, and permitted uses. A structural model $m$ maps this policy into reported receipts, deductions, audits, off-platform activity, and corrected collections. Define $R(d;m)$ as net valid collections, $C_A(d;m)$ as administration cost, $C_B(d;m)$ as business/individual compliance cost, and $L(d;m)$ as expected monetized privacy and security loss under stated preferences and breach probabilities. Fix a welfare population, weights $w_i\geq0$, and monetary public-fund value $\lambda\geq0$. Let $EV_i(d;m)$ denote all other private welfare changes, excluding separately counted costs. The feasible set imposes the actual legal reporting authority and technical capacity of the jurisdiction studied; no current reporting mandate is inferred from a general OECD webpage.

\subsection*{Formal research question}
\emph{Editorial formulation: the mathematical target below is not a theorem quoted from the references.}

Characterize the marginal and discrete welfare comparisons
\[
 \Delta W_m=\mathbb E_m[w_i\Delta EV_i]+\lambda\Delta R-\Delta C_A-\Delta C_B-\Delta L
\]
for reporting expansions, with each term measured in the same monetary welfare unit and no double counting. Determine what audit-validation samples and reform variation identify the corrected-revenue effect independently of receipt reclassification, expense inflation, and movement outside observed platforms. If $L$ is not point identified, derive the range of its marginal values at which a policy ranking reverses. The goal is a feasible reporting frontier or $\varepsilon$-optimal scope conditional on those valuations, not maximization of the quantity of information collected.

\subsection*{Interpretation and scope}
A rise in third-party-reported sales can be offset by newly reported expenses or relocation of transactions, so reported gross receipts alone do not measure additional tax. Platform data miss agents who leave the platform after regulation; linking independent tax and audit samples is needed to define the full affected population. Privacy is a preference- and institution-dependent loss, not a freely chosen residual guaranteeing a desired result. Incidence matters if reporting burdens fall on small firms while additional receipts come from high-income taxpayers; fixed weights make that trade-off explicit. Network effects can propagate through suppliers and customers, requiring exposure definitions or a model of interference. The extension adds cost, privacy, displacement, and marginal policy design to existing enforceability results.

\subsection*{Source audit and status}
The three original academic references are verified, with [S3] resolved to the US credit-card reporting study. The atlas's transparency-network webpage is contextual and does not specify legal obligations. The cited studies support reporting mechanisms, not a general conclusion that maximal reporting is optimal or still an unsolved theorem.

\subsection*{References}
\begin{enumerate}[label={[S\arabic*]},leftmargin=*,itemsep=0.6em]
\item Henrik Jacobsen Kleven, Martin B. Knudsen, Claus Thustrup Kreiner, Søren Pedersen, and Emmanuel Saez (2011). Unwilling or Unable to Cheat? Evidence from a Tax Audit Experiment in Denmark. Econometrica 79(3), 651--692. DOI: 10.3982/ECTA9113.
\url{https://onlinelibrary.wiley.com/doi/10.3982/ECTA9113}
\newline\emph{Locator and support limits:} Abstract and audit experiment: evasion and third-party information; no general marginal privacy-cost schedule.
\item Dina Pomeranz (2015). No Taxation without Information: Deterrence and Self-Enforcement in the Value Added Tax. American Economic Review 105(8), 2539--2569. DOI: 10.1257/aer.20130393.
\url{https://www.aeaweb.org/articles?id=10.1257/aer.20130393}
\newline\emph{Locator and support limits:} Abstract and experiments: Chilean VAT information and supply-chain enforcement; no comprehensive privacy-cost estimate.
\item Joel Slemrod, Brett Collins, Jeffrey L. Hoopes, Daniel Reck, and Michael Sebastiani (2017). Does Credit-Card Information Reporting Improve Small-Business Tax Compliance? Journal of Public Economics 149, 1--19. DOI: 10.1016/j.jpubeco.2017.02.010.
\url{https://www.sciencedirect.com/science/article/pii/S0047272717300233}
\newline\emph{Locator and support limits:} Abstract: US Form 1099-K reporting and sole proprietors; reported expenses and receipts need joint analysis.
\end{enumerate}

\subsection*{Common conventions: Source mathematical conventions}

Unless an entry states otherwise, agent and firm sets are finite. State and action spaces are measurable, strategies and outcome maps are measurable, probability laws are specified, and expectations exist for the targets under discussion. Infinite-horizon discount factors lie in $(0,1)$ and discounted objectives are finite. Norms, tolerances and complexity input sizes must be specified for computational comparisons. Constants or primitives that appear locally are inputs, not empirical facts asserted by this catalogue.

\textbf{Identification.} A statistical model $m\in\mathcal M$ implies an observable law $P_m$ and target $\theta(m)$. At law $P$, its identified set is
\[
\Theta(P;\mathcal M)=\{\theta(m):m\in\mathcal M,\ P_m=P\}.
\]
Point identification holds when this set is a singleton, or equivalently when $P_{m_1}=P_{m_2}$ implies $\theta(m_1)=\theta(m_2)$ for all compatible models. Sharp bounds mean bounds attained, or approached when the boundary is not attained, by compatible models. Writing this set is a definition, not a solution: the substantive task is to establish informative restrictions and derive or estimate the set. An unrestricted-confounding model may give only trivial bounds.

\textbf{Inference.} Identification, estimability and valid uncertainty quantification are separate questions. Fix $\alpha\in(0,1)$. A confidence procedure $C_n$ over a declared law class $\mathcal P$ has asymptotic uniform coverage at level $1-\alpha$ when
\[
\liminf_{n\to\infty}\inf_{P\in\mathcal P}\Pr_P\{\theta(P)\in C_n\}\geq1-\alpha.
\]
For set-identified targets the coverage event must be defined separately, for example coverage of the full identified set. If an entry uses identification-indexed models instead of uniquely indexed laws, the same coverage convention is applied over those models. Pointwise limits do not establish uniform validity, and asymptotic validity does not guarantee finite-sample precision. Sample sizes, dependence structures and weak-identification sequences are part of each application.

\textbf{Counterfactuals.} $Y_i(a)$ is a potential outcome under a specified intervention, including its timing, implementation and versions. It is not $Y_i$ conditional on voluntarily choosing $a$. A full intervention vector permits interference; shorthand scalar treatment notation does not silently impose no interference. Assignment, selection, attrition, anticipation and measurement must be specified. A claim about an instrument's local effect is not automatically a population policy effect.

\textbf{Economic equilibrium.} A counterfactual model includes best responses, constraints, feasibility, market clearing, state transitions and expectations consistent with the stated information. When several equilibria exist, either a declared selector is an input or the target is set-valued. Comparisons across policy regimes must use the stated selection convention. Reduced-form relationships are not presumed invariant after an equilibrium-changing intervention.

\textbf{Optimization and welfare.} Optimization is over the stated feasible set. If attainment is not established, $\sup$ or $\inf$ and $\varepsilon$-optimal choices replace an asserted maximizer or minimizer. For example, with value $V=\sup_{a\in\mathcal A}W(a)<\infty$, an $\varepsilon$-optimal set is $\{a\in\mathcal A:W(a)\geq V-\varepsilon\}$ for $\varepsilon>0$. Benefits, costs, risk and health or environmental outcomes can be aggregated only using declared units and conversion weights. Interpersonal utility weights, the treatment of fiscal transfers and foreign profits, discounting, and the behavioral welfare criterion are explicit inputs. A comparative-static derivative additionally requires existence and appropriate differentiability of the selected counterfactual.

\textbf{Sources.} Local labels [S1], [S2], and so forth restart in every question. Locators identify the relevant abstract, model, section or result at the level actually checked; a bibliographic or abstract-level verification is not represented as inspection of an entire proof. Working papers are distinguished from later publications and changing versions. Source summaries are paraphrased. All URL checks and editorial formulations refer to the audit date stated above.
% END ECONOMICS STATEMENT
\end{document}
